\begin{table*}[htbp]
\centering
\caption{Characteristics of Respiratory Sound Types}
\label{tab:resp_chars}
\small
\resizebox{\textwidth}{!}{ % Resize table to fit page width
\renewcommand{\arraystretch}{1.8}
\begin{tabular}{>{\centering\arraybackslash}m{2.5cm}>{\centering\arraybackslash}m{1.5cm}>{\centering\arraybackslash}m{1.5cm}>{\centering\arraybackslash}m{2.2cm}>{\centering\arraybackslash}m{1.4cm}>{\centering\arraybackslash}m{1.2cm}>{\centering\arraybackslash}m{4.5cm}}
\toprule
\textbf{Disease} & \textbf{Frequency (Hz)} & \textbf{Duration (ms)} & \textbf{Timing in Respiratory Cycle} & \textbf{Continuity} & \textbf{Type} & \textbf{Additional Characteristics} \\ \midrule
\textbf{Pneumonia} & 600-650 & ± 5 & Late inspiration, 
early expiration & Yes & Fine Crackle & High-pitched, crackling sounds like hair crackling near a fire; may change with position \\ 
\textbf{Lung Fibrosis} & 200-350 & ± 15 & Late inspiration & No & Coarse Crackle & Low-pitched, crackling sound like cellophane crinkling; persistent \\ 
\textbf{Chronic Bronchitis, COPD} & 100-500 & 80-100 & Throughout expiration & Yes & Wheeze & High-pitched, musical whistling sounds may vary in intensity and pitch \\ 
\textbf{Bronchiectasis} & 50-200 & 250 & Throughout expiration & Yes & Rhonchus & Low-pitched, snoring, or bubbling sounds, often heard in larger airways \\ 
\textbf{Epiglottitis, Croup} & 200 & 200 & Inspiration & Yes & Stridor & The high-pitched, harsh, crowing sound indicates airway obstruction \\ 
\textbf{Laryngeal Edema} & 200 & 200 & Inspiration & Yes & Squawk & A high-pitched, squeaking sound suggests the narrowing of the upper airway \\ 
\textbf{Pneumonia, Hypersensitivity} & 200 & 15 & Late inspiration & Yes & Pleural Rub & Scratching or grating sound indicates inflammation of the pleura (lung lining) \\ 
\textbf{Lung Tumor} & 200 & 250 & Variable & Yes & Gasp & A harsh, gasping sound may indicate severe respiratory distress \\ \bottomrule
\end{tabular}
}
\end{table*}
